% NARRATIVE: Establish the technical foundation. Reader needs to understand TRM
% and ACT before we can explain halt-first selection and oracle-first training.
\section{Background}

\subsection{Tiny Recursive Models}

Tiny Recursive Models (TRM)~\citep{jolicoeur2025trm} achieve strong performance
on constraint satisfaction problems using a remarkably simple architecture: a
small neural network that iteratively refines its predictions through recursive
application:
%
\begin{algorithm}[h]
\caption{TRM with ACT (Inference)}
\label{alg:trm}
\begin{algorithmic}[1]
  \FOR{$t = 0, \ldots, n_\text{ACT}{-}1$ \LCOMMENT{$n_\text{ACT}{=}16$}}
    \FOR{$i = 0, \ldots, n_\text{H}{-}1$ \LCOMMENT{$n_\text{H}{=}6$}}
      \STATE $z_L, z_H \gets \operatorname{stop}(z_L), \operatorname{stop}(z_H)$
      \FOR{$j = 0, \ldots, n_\text{L}{-}1$ \LCOMMENT{$n_\text{L}{=}9$}}
        \STATE $z_L \gets f(\operatorname{embed}(x;\theta) + z_L + z_H;\theta)$
      \ENDFOR
      \STATE $z_H \gets f(z_L + z_H;\theta)$
    \ENDFOR
    \STATE $y,q \gets h_\text{puzzle}(z_H;\theta), h_\text{halt}(z_H;\theta)$
    \IF{$q > 0$}
      \STATE \textbf{break}
    \ENDIF
  \ENDFOR
\end{algorithmic}
\end{algorithm}
%
The model maintains two latent states: $z_L$ for ``low-level'' details (updated
$n_\text{L}$ times per H-cycle, i.e., one outer loop iteration) and $z_H$ for
``high-level'' reasoning (updated once per H-cycle). The $\operatorname{stop}(\cdot)$ operator detaches
its argument from the computation graph, preventing gradient flow between
H-cycles and enabling stable training.

On Sudoku-Extreme (17-clue puzzles, the minimum for uniqueness), a 7M parameter
TRM achieves ${\sim}$86\% puzzle accuracy---substantially outperforming larger
language models while using a fraction of the parameters.

\subsection{Adaptive Computation Time}

Rather than fixing the number of iterations $T$, TRM can learn \emph{when to
stop} via Adaptive Computation Time (ACT)~\citep{graves2016act}. The model
outputs a halting probability $q_\text{halt}\defeq\sigmoid(q)$ at each step,
trained to predict whether the current output is correct. When $q_\text{halt} >
0.5$, the model ``commits'' to its answer.

This creates an implicit tradeoff: easy puzzles should halt quickly (high
confidence after few iterations), while hard puzzles may require more
computation. In practice, a well-trained TRM solves most Sudoku puzzles in 1--3
steps, with a long tail requiring up to 16 iterations.
